% 第二章 相关技术

\chapter{相关技术}

本章围绕FreeFlow系统实现所依赖的关键技术展开说明。结合链上音乐发行、访问控制、资源存储与桌面端播放等业务需求，系统采用桌面端应用、链下服务、智能合约与内容寻址存储协同的技术方案。以下从应用架构、身份认证、资源存储、链上合约与播放支撑五个方面介绍相关技术基础。

\section{桌面端与链下服务技术基础}

FreeFlow以Electron与React构建桌面端，以Node.js、Express和Prisma构建链下服务。该组合能够覆盖桌面应用交互、本地能力调用、HTTP接口组织与关系型数据管理等核心需求。

Electron提供跨平台桌面应用运行环境，并通过主进程与渲染进程分工承载系统功能\cite{electronDocs,electronIPC}。在本系统中，主进程负责窗口管理、本地数据库访问、下载与缓存控制、音频代理等本地能力；渲染进程负责创作者工作台、资源检索、详情展示、评论与播放控制等界面交互。React用于组织前端组件与状态流，使播放器、作品列表、评论区和发行工作台等界面模块具有较好的可维护性\cite{reactDocs}。

链下服务采用Node.js与Express构建。Node.js适合处理上传、流式转发和链上交互代理等I/O密集型任务，Express则便于组织认证、资源、购买、评论和异常处理等业务接口\cite{nodeDocs,expressDocs}。数据持久化方面，系统使用PostgreSQL保存用户、钱包身份、发行记录、上传记录、购买记录与评论数据，并通过Prisma完成模式定义、关系映射与迁移管理\cite{prismaDocs}。该技术组合为后续的资源索引、分页查询、业务状态维护和系统联动提供了稳定的数据支撑。

\section{SIWE身份认证机制}

系统身份认证采用SIWE（Sign-In with Ethereum）方案，其消息格式由EIP-4361定义\cite{eip4361}。该机制以钱包签名证明用户对链上地址的控制权，能够将系统身份与链上账户建立统一关联，便于后续购买、发布和访问校验等流程复用同一身份基础。

在具体实现中，后端首先生成一次性随机数nonce，并构造包含域名、地址、URI、链ID、签发时间等字段的标准化消息；客户端调起钱包完成签名后，将签名结果与原始消息提交至后端；后端再对消息格式、nonce有效性、链ID与签名恢复地址进行校验\cite{eip4361}。验证通过后，系统创建或绑定用户实体、钱包身份与会话记录，并将后续业务操作统一关联至该地址身份。

相较于传统账户口令体系，SIWE减少了独立密码存储与找回流程带来的维护负担，也增强了桌面端身份体系与链上访问控制之间的一致性，因此适合本系统的应用场景。

\section{IPFS内容寻址存储与混合架构}

链上音乐系统同时涉及媒体资源存储、可验证引用和高频业务协同。基于这一特点，FreeFlow将区块链、IPFS与链下服务组合为混合架构，以兼顾可信记录、资源可定位性与系统运行效率。

IPFS采用内容寻址机制，以CID标识文件内容。当资源内容发生变化时，其对应CID也随之变化，因此适合用于音频文件、封面图和metadata等静态对象的发布与引用\cite{ipfsDocs}。本系统使用Pinata完成上传与持久化管理，在创作者发布过程中保存音频、封面和metadata对象，并获得相应CID或网关地址\cite{pinataDocs}。考虑到IPFS本身主要负责寻址与分发，资源的长期可用性仍依赖持续保存机制，系统借助Pinata的持久化服务降低节点维护与pinning管理的复杂度\cite{ipfsPinning}。

在系统层面，作品发布结果、访问凭证、交易记录和收益规则等关键状态保存于链上；评论、资源索引、登录会话、上传状态与检索支撑数据保存于链下。已有研究表明，链上、链下与混合模式在成本、吞吐、可扩展性与可信性之间存在明确权衡，混合方案在多媒体应用中具有较高的工程可行性\cite{Chen2021,Rikken2019,Alawadi2024,Friedl2024,Eren2025}。对于音乐场景而言，媒体文件体积较大，评论与检索数据更新频繁，若全部依赖链上处理，将显著增加存储与交互成本；将高可信状态与高频服务分层组织，更符合系统当前的规模与实现目标。

\section{智能合约与ERC-1155访问凭证}

链上业务逻辑采用Solidity实现\cite{solidityDocs}。系统围绕作品发布、访问授权与收益分配设计合约结构，其中访问控制以ERC-1155多代币标准为基础\cite{eip1155}。

ERC-1155允许在同一合约内管理多种代币类型，适合将每首作品映射为独立的\texttt{tokenId}访问凭证。FreeFlow中，用户完成购买后获得对应\texttt{tokenId}，系统据此进行访问状态校验。由于该凭证在业务语义上承载作品访问资格，合约实现中需要对普通转移行为加以限制，以维持访问控制与购买关系的一致性。

系统合约主要包括\texttt{MusicAccess1155}、\texttt{PlatformHub}和\texttt{RoyaltySplitter}。\texttt{MusicAccess1155}负责访问凭证的创建与校验；\texttt{PlatformHub}负责作品发布、销售参数维护、购买处理与访问判断；\texttt{RoyaltySplitter}负责收益分账与提取。通过上述合约协作，系统可以完成作品发布、凭证发放和收益流转等核心链上业务。为提高标准实现的可靠性，开发过程中复用了OpenZeppelin提供的合约组件与安全工具库\cite{openzeppelinContracts}。

\section{音频代理、Range请求与本地缓存}

在线播放场景对随机访问、进度跳转与重复播放效率具有较高要求。为提升桌面端的播放稳定性，系统在本地集成音频代理服务，并将本地文件、缓存文件与远程音频统一纳入同一访问入口。

播放器通过本地代理接口请求音频数据，代理层根据资源状态决定直接读取本地文件、返回缓存内容，或向远程地址发起转发请求。该方式便于集中处理资源定位、缓存命中、错误恢复与播放链路控制。对于拖动进度条、断点续播等操作，系统依赖HTTP Range机制完成按字节区间读取。根据HTTP语义规范，客户端可通过\texttt{Range}请求头指定目标区间，服务端在支持时返回\texttt{206 Partial Content}响应\cite{rfc9110}。这一机制有助于降低重新加载完整音频文件所带来的时间开销。

实现层面，Node.js提供的HTTP与流式处理能力适合音频转发和缓存写入场景\cite{nodeHttp,nodeStream}。系统对普通播放请求采用边传输边缓存策略，对范围请求优先透传上游分段响应，从而在保证跳转实时性的同时兼顾后续复用效率。该设计为链上音乐播放、试听与下载管理提供了必要的桌面端支撑。
